% ===== sec/11_conclusion.tex =====
\section{Conclusion}
\label{sec:conclusion}

Koohestani et al.~\cite{koohestani2025automata} asked whether
agents are just automata.  We use that question to motivate a
narrower operational claim: agents can be analysed through the
trace languages of the operations they emit, and regular verifier
languages can validate those traces with explicit complexity
bounds.  Within this theory:
\begin{itemize}
\item Agents are 7-tuples
$(\Sigma_\mathcal{A},Q_\mathcal{A},q_0^\mathcal{A},X_\mathcal{A},
\delta_\mathcal{A},F_\mathcal{A},h_\mathcal{A})$ over a fixed
operation alphabet
(Section~\ref{sec:trace}, Definition~\ref{def:agent}).
\item Trace languages classify into the standard
Chomsky hierarchy (Section~\ref{sec:chomsky}), with the
indexed-grammar zone of Aho~\cite{aho1968indexed}
appearing naturally between Type-2 and Type-1.
\item The producer--consumer construction of
Dijkstra~\cite{dijkstra1968cooperating}, instantiated as a
idealised Goal-Driven producer constrained by a Type-3 Data-Driven
consumer, suffices: their intersection is closed at the
producer's class and decidable in linear time at the consumer's
DFA (Theorem~\ref{thm:closure},
Corollary~\ref{cor:budget}).
\item The Type-3 consumer is realised as the deterministic
Aho--Corasick~\cite{aho1975corasick} multi-pattern matcher
compiled from the unit-test catalogue.
\item Type-2, indexed-zone, and Type-1 sub-agents are not
designed; they are projected from the binding
(Theorem~\ref{thm:derived}).
\item A Type-3 agent may itself act as a producer,
backward-inferring constrained changes to the tape, provided the
inference terminates in a decidable verdict; otherwise the
halting problem returns and the verifier-side budget collapses
(Section~\ref{sec:type3-producer}).
\item The undecidability shadow on the Type-0 envelope is the
shadow Robinson and the MRDP
theorem~\cite{robinson1949definability,davis1961decision} cast
on every recursively enumerable set.
\end{itemize}

The architecture is not metaphorical.  Instantiated on the
\textsc{Spaceship Titanic} task drawn from the public
\textsc{MLE-bench} benchmark, it produced a leaderboard placement
of $710$ out of $2330$ teams ($\approx\!30.5\%$, top tertile),
representing a $2$--$5\%$ absolute placement improvement over
prior published agent baselines.  The Goal-Driven producer had
no a priori knowledge of the goal label; the Data-Driven
consumer had no LLM component.  The result was obtained by binding the two principal roles and
enumerating a finite validated trace for the task instance.

Section~\ref{sec:verification_ablation_study} supplies the complementary
claim: in a controlled suite, $\mathcal{G}\,\Vert\,\mathcal{V}_{\mathrm{DFA}}$
raises mean sandbox score from $35.2\%$ to $57.8\%$, drives undetected
errors to zero, and outperforms an LLM judge on accuracy at negligible
verifier cost.  The TLV runtime (\path{framework/}, \path{src/}) makes the
verifier, projection engine, and language analyzer executable artifacts,
not paper-only notation.  All reproducible code, including the
start-language CSV, the recursively enumerated trace CSV, sandbox CSVs
under \path{data/}, and the notebooks that produced the submission, is
published at\\
\url{https://www.kaggle.com/scottweeden/competitions}.

The realised trace is the empirical witness to the theory.  Its
thirteenth column, \texttt{kaggle\_trainer}, is the emergent
agent that the binding forced into existence: the start language
had twelve sub-agents and no trainer; the realised language has
thirteen sub-agents and a 36-step training trace.  The expansion
is the operational form of the Type-0 recursive enumeration; the
laptop's finite tape is the operational form of the Type-1 bound
that brings the enumeration to a halt.

\subsection{Future work}

Four directions are immediate.  First, the sandbox ablation should
be repeated on additional live \textsc{MLE-bench} tasks with real API
judges to test external validity beyond
Section~\ref{sec:verification_ablation_study}.  Second, the start language should
be examined as a design object in its own right: which start
languages produce stable bindings, and which fail to converge?
Third, the trace-language identity criterion of
Definition~\ref{def:traceeq} invites a formal study of agent
equivalence: when are two agents with different LLM backbones
trace-equivalent on a given task?  Fourth, suggested by the
Section~\ref{sec:type3-producer} discussion, is to characterise
the class of regular keyword-update rules that preserve
decidability under a Type-3 producer; we conjecture this class
is exactly the regular extensions of the start-language alphabet.

\subsection{Closing remark}

Turing wrote in 1936 that ``the justification [for the
definitions] lies in the fact that the human memory is
necessarily limited''~\cite{turing1936computable}.  Our laptop's
memory is necessarily limited.  Our agent's tape is necessarily
limited.  Within those limits, the trace-validation claims are exact;
outside them, the Type-0 abstraction is an idealisation.

A final note.  This paper began as a question in an email about
arithmetic grammars and grew under conditions in which other
deliverables were not available to the first author --- a
sequence we record honestly in the Acknowledgments.  Papers are
sometimes written in a straight line and sometimes in a crooked
one, and the architecture proposed here was, in a real sense,
shaped by both: the formal narrowness of the theory, and the
narrow window in which the theory had to be made to do real
work.  We close with the narrower claim that an agent can be studied
through its trace language, and that this viewpoint gives useful
producer--consumer tools for language-constrained trace validation.
